\section{Discussion and Limitations}
\label{section:discussion}

\paragraph{Why the AR pitch failure was misread.} The headline lesson is
methodological: an architecture conclusion (``a from-scratch ResNet cannot learn
pitch'') was really a statement about the \emph{decoder's alignment}, not the
encoder. Because our four-head design measures pitch separately, and because the
CTC variant reuses the identical encoder features, we could falsify the encoder
explanation directly. The decoupled-stream contract was what made the negative
result interpretable rather than merely discouraging.

\paragraph{Revisiting the alignment bottleneck formally.} The same argument can be
stated at the level of the decoding objective. Write
$\mathbf{H}=(\mathbf{h}_1,\dots,\mathbf{h}_T)$ for the column-pooled ResNet
features. The AR decoder factorises the target likelihood autoregressively,
\begin{equation*}
  P(\mathbf{y}\mid\mathbf{H})=\prod_{u} P\!\left(y_u \mid y_{<u},\,\mathbf{H}\right),
\end{equation*}
conditioning each step on a cross-attention context
$\mathbf{c}_u=\sum_{t}\alpha_{u,t}\,\mathbf{h}_t$. Because column pooling makes the
features translation-equivariant---discarding the \emph{absolute} vertical position
that \emph{is} pitch on a staff---and the attention mass $\alpha_{u,t}$ concentrates
on a narrow temporal window per symbol, the decoder is starved of the global
staff-line context needed to infer absolute pitch. CTC instead marginalises over
every valid alignment path $\boldsymbol{\pi}$ under the many-to-one collapse
$\mathcal{B}$,
\begin{equation*}
  P(\mathbf{y}\mid\mathbf{H})=\!\!\!\sum_{\boldsymbol{\pi}\in\mathcal{B}^{-1}(\mathbf{y})}\ \prod_{t} P\!\left(\pi_t \mid \mathbf{H}\right),
\end{equation*}
reading from BiLSTM-contextualised features
$\mathbf{h}'_t=\mathrm{BiLSTM}(\mathbf{H})_t$. The marginalisation dissolves the
rigid one-column-per-symbol commitment, while the bidirectional context lets each
column aggregate the staff-line geometry of its neighbours into the vertical
reference frame that column pooling removed---recovering pitch in full, and without
any explicit 2-D positional embedding. This is why the \emph{identical} encoder
features that look pitch-blind under AR (\S\ref{sec:exp:arch}) become perfectly
legible under CTC.

\paragraph{Pretraining is not a free low-data prior.} The sample-efficiency gap
(\S\ref{sec:exp:scaling}) cautions against assuming ImageNet-pretrained encoders
help most when data is scarce. For a task whose key cue---absolute vertical
position---is orthogonal to natural-image semantics, the prior contributes little,
and the inductive bias of the decoder (CTC's alignment flexibility) dominates
sample efficiency.

\paragraph{Limitations.} (i)~\emph{Scope.} We handle only printed, monophonic
notation; chords, slurs, multi-voice scores, and handwriting are out of scope, and
real-world utility on polyphonic or hand-written music remains untested.
(ii)~\emph{No external baselines.} Our comparisons are internal (ViT-AR vs
ResNet-AR vs CRNN+CTC). We did not run the pipeline OMR baseline (Oemer) or a
pretrained Vision-Encoder-Decoder OCR baseline (TrOCR) end-to-end on this task; a
head-to-head against those systems is future work. (iii)~\emph{Statistical power.}
The scaling sweep uses a single seed per cell, so we report trends, not confidence
intervals; the ViT-AR data threshold in particular may move with seed and
data-loader settings, though the order-of-magnitude gap to CRNN+CTC is stable.
(iv)~\emph{Sim-to-real ceiling.} Even after fine-tuning, real-photo pitch accuracy
is $64.5\%$---useful but far from the clean $99.8\%$; the dominant residual is
camera-angle vertical mis-registration, which the height-preserving CRNN+CTC
encoder is well-positioned to attack but which we have not yet combined with the
photo augmentation. (v)~\emph{Deployment gap.} The best model (CRNN+CTC) is not
yet the deployed one: its CTC decoder needs the type-anchored collapse wired into
the inference path, and its recurrent weights require a serialization fix to
survive checkpoint reloading. The deployed demo therefore ships the ViT-AR model.

\paragraph{Future work.} The natural next step is to bring CRNN+CTC---which is both
more accurate and far more sample-efficient---through deployment, then fine-tune
\emph{it} with the photo augmentation, since its preserved vertical resolution
directly targets the sim-to-real bottleneck. Adding the external Oemer/TrOCR
baselines would situate our system against the broader OMR literature, and
extending the generator toward polyphony would test how far the four-stream
contract scales.
